\largerpage
\sscp{Tree kangaroos (\it{Dendrolagus}) and wallabies}{C-06-03}
Data taken from word lists is usually unclear as to which kangaroo-like
animal the elicited term designates.
We follow the glosses in the sources,
though note that the same term is sometimes is glossed by one
source as ‘kangaroo/wallaby’ and sometimes by another source as ‘tree kangaroo/wallaby’.
Terms glossed ‘tree wallaby/kangaroo’ are given before terms glossed ‘kangaroo/wallaby’.
\mbox{}\\

\noindent{\wg\st{2pt}\begin{tabularx}{\linewidth}
{lllll>{\raggedright\arraybackslash}X}
\lsptoprule
%Language	&	ISO	&	Term	&	Gloss	&	Etymon	&	Source	\\	\midrule
%\endfirsthead
Language	&	ISO	&	Term	&	Gloss	&	Etymon	&	Source	\\	\midrule
%\endhead
Kowiai-1		&	kwh	&	ɸanai			&	tree wallaby	&	\it{misc.}		&	\cite{WalkerWalker1991} (Mly. \it{lao-lao)}	\\
Kowiai-2		&	kwh	&	waŋguria	&	tree wallaby	&	**waked				&	SV \citeyear[263]{SmitsVoorhoeve1992b}	\\
Uruangnirin	&	urn	&	taˈker		&	tree wallaby	&								&	\cite{Visser2023}	\\
Onin				&	oni	&	kisor;		&	tree wallaby;	&	{\#}kisor(o);	&	SV \citeyear[263]{SmitsVoorhoeve1992b};	\\	\hhline{~}
						&			&	kerá			&	tree wallaby	&	\it{misc.}		&	\cite{Donohue2010}	\\
Sekar				&	skz	&	tamtik		&	tree wallaby	&	\it{misc.}		&	SV \citeyear[263]{SmitsVoorhoeve1992b}	\\
Arguni			&	agf	&	oyambo		&	tree wallaby	&	\it{misc.}		&	SV \citeyear[263]{SmitsVoorhoeve1992b}	\\
Erokwanas		&	erw	&	wayambu		&	tree kangaroo	&	\it{misc.}		&	\cite{Donohue2010}	\\
Irarutu			&	irh	&	waire;		&	tree wallaby;	&	**waked;			&	SV \citeyear[263]{SmitsVoorhoeve1992b};	\\	\hhline{~}
						&			&	amoro			&	tree wallaby	&	**amoR				&	SV \citeyear[263]{SmitsVoorhoeve1992b}	\\
\lspbottomrule
\end{tabularx}}

\noindent{\wg\st{2pt}\begin{xltabular}{\linewidth}
{lllll>{\raggedright\arraybackslash}X}
\lsptoprule
%Language	&	ISO	&	Term	&	Gloss	&	etymon	&	source	\\	\midrule
%\endfirsthead
Language	&	ISO	&	Term	&	Gloss	&	etymon	&	source	\\	\midrule	\hhline{~}
\endhead
Kei\footnotemark&	kei	&	laʔe;	&	kangaroo; 		&	\it{misc.}	&	\cite{Travis2019}	\\
							&			&	aha;		&	wallaby; 			&	\it{misc.}	&	(Mly. \it{kangguru pohon} `tree kangaroo')	\\
							&			&	ha			&	wallaby				&	\it{misc.}	&	(Mly. \it{sm kangguru yg agak kecil} `k.o. kangaroo which is rather small')	\\
Uruangnirin		&	urn	&	taˈpar	&	k.o. kangaroo	&	\it{misc.}	&	\cite{Visser2023} (Mly. \it{lalau})	\\
Onin					&	oni	&	wekar		&	kangaroo			&	**waked			&	SV \citeyear[241]{SmitsVoorhoeve1992b}	\\
Sekar					&	skz	&	wekar;	&	kangaroo;			&	**waked;		&	SV \citeyear[241]{SmitsVoorhoeve1992b};	\\	\hhline{~}
							&			&	lawlaw	&	kangaroo			&	\it{misc.}	&	\cite{Donohue2010} (Mly. \it{laolao/saham})	\\
Arguni				&	agf	&	utae;		&	kangaroo;			&	\it{misc.}	&	SV \citeyear[241]{SmitsVoorhoeve1992b};	\\	\hhline{~}
							&			&	ayámbu	&	kangaroo			&	\it{misc.}	&	\cite{Narfafan2011}	\\
Irarutu				&	irh	&	wasio		&	kangaroo			&	\it{misc.}	&	SV \citeyear[241]{SmitsVoorhoeve1992b}	\\	\hhline{~}
							&			&	amore		&	kangaroo			&	**amoR			&	SV \citeyear[241]{SmitsVoorhoeve1992b}	\\
Kuri					&	nbn	&	kabˈey	&	kangaroo			&	\it{misc.}	&	SV \citeyear[241]{SmitsVoorhoeve1992b}	\\
\lspbottomrule
\end{xltabular}}
\footnotetext{
		Only one species of macropod,
		\it{Thylogale brunii} ‘Dusky Pademelon’,
		is known to occur in the Kei Islands.
		How exactly the different Kei terms for `kangaroo'
		and `wallaby' relate to this macropod is unclear.}